%! TEX root = ../note-dqm-transform.tex

\section{Differentiability in quadratic mean}
\label{sec--dqm-def}

Let \((\mathbf{X}, \mathcal{B})\) be a measurable space.
We use linear functional notation for integrals: for a measure \(\mu\) on
\((\mathbf{X}, \mathcal{B})\), \(\mu [f] := \int f (x) \; \mu (\mathrm{d} x)\).
Furthermore let \(d \in \mathbb{N}\) and \(\Theta \subseteq \R^{d}\); denote the
Euclidean norm on \(\R^{d}\) by \(|\cdot|\), i.e. \(|x|^{2} = \sum_{j = 1}^{d}
x_{j}^{2}\) for any \(x \in \R^{d}\).
Throughout, we let \(\mu\) be a \(\sigma\)-finite measure on \((\mathbf{X},
\mathcal{B})\) and let \(\Theta\) be a non-empty subset of \(\R^{d}\).

Let \(\mathbf{Q} = \left\{ Q_{\theta} : \theta \in \Theta \right\}\) be a family
of probability measures dominated by \(\mu\), a \(\sigma\)-finite positive
measure.
These are possible distributions for a random element \(X\) with values in
\(\mathbf{X}\).
Denote the \(\mu\)-densities by
\begin{equation}
  q_{\tau} (\cdot) := q (\cdot; \tau) := \frac{\mathrm{d}
  Q_{\tau}}{\mathrm{d} \mu}.
  \label{eqn--q-density-and-sqrt}
\end{equation}
Since \(q_{\theta}\) is a \(\mu\)-density for a probability measure, its square
root inhabits \(\L_{2} (\mu)\).
Differentiability in quadratic mean (DQM) is essentially Fr\'echet
differentiability of the root-density map \(\theta \mapsto \sqrt{q_{\theta}}\),
with some additional restrictions on the structure of the derivative.
We first revisit the definition of Fr\'echet differentiation in
\(\L_{r} (\mu)\) for generic \(r \in [1, \infty)\) in
\Cref{def--frechet-diff-relative}.

\begin{definition}[Relative Fr\'echet differentiability in \(\L_{r}
(\mu)\)]
\label{def--frechet-diff-relative}
Let \(\theta_{0} \in \Theta\) and \(r \in [1, \infty)\).
\(g : \Theta \to \L_{r} (\mu)\) is Fr\'echet differentiable at
\(\theta_{0}\) relative to \(\theta_{0}\) in \(\L_{r} (\mu)\) if there
is a \(d\)-vector of \(\L_{r} (\mu)\) functions \(\dot{g}_{\theta_{0}}
(\cdot) = \dot{g} \left( \cdot; \theta_{0} \right)\) such that
\begin{equation}
  \begin{gathered}
    \mu \left[ \left| \rho_{\theta, \theta_{0}} \right|^{r} \right] = o \left(
    \left| \theta - \theta_{0} \right|^{r} \right) \text{ as } \theta \to
    \theta_{0} \text{ while } \theta \in \Theta, \\
    \text{where} \quad
    \rho_{\theta, \theta_{0}} = g_{\theta} - g_{\theta_{0}} - \left(
    \theta - \theta_{0} \right)^{\prime} \dot{g}_{\theta_{0}}.
  \end{gathered}
  \label{eqn--diff-def-remainder}
\end{equation}
Say that \(g\) is Fr\'echet differentiable relative to \(\theta_{0}\) in
\(\L_{r} (\mu)\) if \eqref{eqn--diff-def-remainder} holds at every
\(\theta_{0} \in \theta_{0}\).
\end{definition}

Henceforth, we shall refer to relative Fr\'echet differentiability in
\(\L_{r}\) as simply ``\(\L_{r}\)-differentiability'' for
brevity.
Some useful \(\L_{1}\)-differentiability and contiguity results for
densities arise solely from \(\L_{2}\)-differentiability of the
root-density maps.
At interior points of \(\Theta\), we get two additional benefits for
non-negative maps:
faster ``contiguity rates'' and nullity of the derivative
\(\dot{g}_{\theta_{0}}\) on \(Q_{\theta_{0}}\)-null sets.

\begin{lemma}
\label{lem--hellinger-diff-properties}
Let \((\mathbf{X}, \mathcal{B}, \mu)\) be a measure space and let \(\theta
\mapsto g_{\theta}\) be a map from \(\Theta\) into \(\L_{2} (\mu)\)
that is \(\L_{2} (\mu)\)-differentiable at
\(\theta_{0} \in \Theta\) with derivative \(\dot{g}_{\theta_{0}}\).
\begin{enumerate}[label=(\alph*)]
  \item \label{lem--hellinger-diff-properties-q-diff-in-L1}
    The map \(\theta \mapsto q_{\theta} := g_{\theta}^{2}\) is \(\L_{1}
    (\mu)\)-differentiable at \(\theta_{0}\) with derivative \(\dot{q}_{\theta_{0}} = 2
    g_{\theta_{0}} \dot{g}_{\theta_{0}}\).
  \item \label{lem--hellinger-diff-properties-q-singularity-result}
    If \(g_{\theta} \geq 0\) \(\mu\)-a.e for every \(\theta \in \Theta\), then
    \begin{equation}
      \mu \left[ g_{\theta}^{2} \cdot \mathbf{1} \left\{ g_{\theta_{0}} = 0
      \right\} \right] = O \left( |\theta - \theta_{0}|^{2} \right).
      \label{eqn--hellinger-diff-properties-q-singularity-result-1}
    \end{equation}
    Furthermore, if \(\theta_{0} \in \mathrm{int} (\Theta)\), then
    \begin{equation}
      \mu \left[ g_{\theta}^{2} \cdot \mathbf{1} \left\{ g_{\theta_{0}} = 0
      \right\} \right] = o \left( |\theta - \theta_{0}|^{2} \right).
      \label{eqn--hellinger-diff-properties-q-singularity-result}
    \end{equation}
  \item \label{lem--hellinger-diff-properties-dotg-singularity-result}
    If \(g_{\theta} \geq 0\) \(\mu\)-a.e for every \(\theta \in \Theta\) and
    \(\theta_{0} \in \mathrm{int} (\Theta)\), then
    \begin{equation}
      \mu \left[ \left| \dot{g}_{\theta_{0}} \right|^{2} \cdot \mathbf{1}
      \left\{ g_{\theta_{0}} = 0 \right\} \right] = 0.
      \label{eqn--hellinger-diff-properties-gdot-singularity-result}
    \end{equation}
\end{enumerate}
\end{lemma}

\begin{proof}[Proof of \Cref{lem--hellinger-diff-properties}]
See \Cref{sec--prf--lem--hellinger-diff-properties}.
\end{proof}

Let \(g_{\theta} = \sqrt{q_{\theta}}\) denote the square roots of densities in
\(\mathbf{Q}\).
In some texts, e.g. \citet{1998bickelEfficientAdaptiveEstimation},
differentiability in quadratic mean is defined as \(\L_{2}\)-differentiability
of \(\theta \mapsto g_{\theta}\) at interior points of \(\Theta\).
Denoting the log-density by \(\ell_{\theta} = \mathbf{1} \left\{
q_{\theta} > 0 \right\} \log q_{\theta} = 2 \mathbf{1} \left\{ g_{\theta} > 0
\right\} \log g_{\theta}\), the score would be defined via the chain rule as
\begin{equation}
  \dot{\ell}_{\theta_{0}} = \frac{2 \cdot \mathbf{1} \left\{ g_{\theta_{0}} > 0
  \right\}}{g_{\theta_{0}}} \dot{g}_{\theta_{0}} \implies \dot{g}_{\theta_{0}}
  = \frac{1}{2} \sqrt{q_{\theta_{0}}} \cdot \dot{\ell}_{\theta_{0}}.
  \label{eqn--score-root-density-deriv}
\end{equation}
The \(\L_{2}\)-differentiability condition would then read:
\begin{equation}
  \mu \left[ \left( \sqrt{q_{\theta}} - \sqrt{q_{\theta_{0}}} - \frac{1}{2}
  \sqrt{q_{\theta_{0}}} \cdot \dot{\ell}_{\theta_{0}}^{\prime} \left( \theta -
  \theta_{0} \right) \right)^{2} \right] = o \left( \left| \theta -
  \theta_{0} \right|^{2} \right).
  \label{eqn--dqm-def-0}
\end{equation}
If we directly posit the existence of a \(\dot{\ell}_{\theta_{0}} \in \L_{2}
\left( Q_{\theta_{0}} \right)\) such that \eqref{eqn--dqm-def-0} holds,
then we get the properties
\ref{lem--hellinger-diff-properties-q-singularity-result} and
\ref{lem--hellinger-diff-properties-dotg-singularity-result} in
\Cref{lem--hellinger-diff-properties} essentially for free.
By doing so, we force
\(\dot{g}_{\theta_{0}} = \frac{1}{2} \sqrt{q_{\theta_{0}}} \cdot
\dot{\ell}_{\theta_{0}}\), whence
\eqref{eqn--hellinger-diff-properties-gdot-singularity-result} in
\Cref{lem--hellinger-diff-properties}
\ref{lem--hellinger-diff-properties-dotg-singularity-result} is trivial.
For \Cref{lem--hellinger-diff-properties}
\ref{lem--hellinger-diff-properties-q-singularity-result}, note that since
\(q_{\theta_{0}} \cdot \ind \left\{ q_{\theta_{0}} > 0 \right\} =
q_{\theta_{0}}\) and \(\ind \left\{ q_{\theta_{0}}
> 0 \right\} \ind \left\{ q_{\theta_{0}} = 0 \right\} = 0\),
\begin{equation*}
  \mu \left[ q_{\theta} \ind \left\{ q_{\theta_{0}} = 0 \right\} \right] \leq
  \mu \left[ \left( \sqrt{q_{\theta}} - \sqrt{q_{\theta_{0}}} - \frac{1}{2}
  \sqrt{q_{\theta_{0}}} \cdot \dot{\ell}_{\theta_{0}}^{\prime} \left( \theta -
  \theta_{0} \right) \right)^{2} \right] = o \left( \left| \theta -
  \theta_{0} \right|^{2} \right).
\end{equation*}

In any case, parts \ref{lem--hellinger-diff-properties-q-singularity-result} and
\ref{lem--hellinger-diff-properties-dotg-singularity-result} of
\Cref{lem--hellinger-diff-properties} carry much of the benefit (but not all) of
\(\L_{2}\)-differentiability of the root-density map \(\theta \mapsto
\sqrt{q}_{\theta}\).
Thus some texts, e.g. \citet{2000lecamAsymptoticsStatisticsBasic}, take these
properties as primitives in when defining the concept differentiability in
quadratic mean.
In fact, even the requirement of a common dominating measure \(\mu\) can be
dispensed with when taking this approach.
Given \(\theta_{0} \in \Theta\), by Lebesgue's Decomposition Theorem we can
write
\begin{equation}
  Q_{\theta} = Q_{\theta, \theta_{0}}^{\perp} + Q_{\theta, \theta_{0}}
  \quad \text{where} \quad
  Q_{\theta, \theta_{0}}^{\perp} \perp Q_{\theta_{0}},
  \quad \text{and} \quad
  Q_{\theta, \theta_{0}} \ll Q_{\theta_{0}},
  \label{eqn--Qtheta-leb-decomp}
\end{equation}
are the parts of \(Q_{\theta}\) that are singular and absolutely continuous
respectively against \(Q_{\theta_{0}}\).
Then part \ref{lem--hellinger-diff-properties-q-singularity-result}
of \Cref{lem--hellinger-diff-properties} translates a requirement that
\(Q_{\theta, \theta_{0}}^{\perp} (\mathbf{X}) = o \left( |\theta -
\theta_{0}|^{2} \right)\).
\Cref{def--dqm} combines this requirement with a version of
\eqref{eqn--dqm-def-0} reformulated for \(\mathrm{d} Q_{\theta,
\theta_{0}} / \mathrm{d} Q_{\theta_{0}}\), noting that the density is equivalent
to \(1\)
\(Q_{\theta_{0}}\)-a.e. when \(\theta = \theta_{0}\).

\begin{definition}[Differentiability in Quadratic Mean]
\label{def--dqm}
\(\mathbf{Q}\) is differentiable in quadratic mean (DQM) at \(\theta_{0} \in
\Theta\) if given the Lebesgue Decomposition in \eqref{eqn--Qtheta-leb-decomp},
the following hold as \(\theta \to \theta_{0}\).
\begin{enumerate}[label=(\roman*)]
  \item \label{def--dqm-sing}
    The singular part of \(Q_{\theta}\) with respect to \(Q_{\theta_{0}}\)
    satisfies
    \(Q_{\theta, \theta_{0}}^{\perp} (\mathbf{X}) = o \left( \left| \theta -
    \theta_{0} \right|^{2} \right)\).
  \item \label{def--dqm-ac-diff}
    Let \(q_{\theta, \theta_{0}} = \mathrm{d} Q_{\theta, \theta_{0}} /
    \mathrm{d} Q_{\theta_{0}}\) be the density of the absolutely continuous part
    of \(Q_{\theta}\) with respect to \(Q_{\theta_{0}}\).
    There exists a \(d\)-vector \(\dot{\ell}_{\theta_{0}} \in \L_{2} \left(
    Q_{\theta_{0}} \right)\) such that
    \begin{equation}
      Q_{\theta_{0}} \left[ \left( \sqrt{q_{\theta, \theta_{0}}} - 1 -
      \frac{1}{2} \dot{\ell}_{\theta_{0}}^{\prime} \left( \theta - \theta_{0}
      \right) \right)^{2} \right] = o \left( \left| \theta - \theta_{0}
      \right|^{2} \right).
      \label{eqn--dqm-def}
    \end{equation}
\end{enumerate}
\end{definition}

For brevity we can simply re-parametrize to \(\theta_{0} = 0\), where we are
denoting the origin in \(\R^{d}\) (at mild risk of abiguity) by \(0\).
Set \(\theta_{0} = 0\) and assume \(\Theta\) contains an open neighborhood of
\(0\).
Then, \(Q_{0}\), \(q_{0}\) and \(\dot{\ell}_{0}\) are the true distribution,
density and score.
Similarly, \(Q_{\theta, 0}^{\perp}\) and \(Q_{\theta, 0}\) are the singular and
absolutely continuous parts of \(Q_{\theta}\) against \(Q_{0}\), and
\(q_{\theta, 0}\) is the density of \(Q_{\theta, 0}\) against \(Q_{0}\).

%%% Local Variables:
%%% mode: LaTeX
%%% TeX-master: "../note-dqm-transform"
%%% End:
